% Tiny contract
% Teams: Suspension
% Requirements: REQ-DES-130 REQ-FUN-020
% Status: [in-progress]

The rover wheels are manufactured as single-piece TPU components to reduce assembly 
complexity and provide controlled elastic compliance on irregular terrain. A chevron-type 
tread is used to improve longitudinal traction and reduce lateral slip on loose soil and 
mixed rocky surfaces typical of the ERC Mars Yard.

To reduce unsprung mass, the wheel body uses a perforated web pattern. The geometry 
was verified by finite-element analysis under a maximum wheel load of \SI{45}{\kilo\gram}, 
with the design target of limiting peak deformation to \SI{5}{\milli\metre}. The current 
wheel revision satisfies this deformation criterion and is used as the baseline for mass, 
traction, and endurance tests.

\begin{figure}[H]
\centering
\begin{minipage}[t]{0.48\linewidth}
	\centering
	\fbox{\parbox[c][4.0cm][c]{0.92\linewidth}{\centering
	Wheel design image\\placeholder}}
\end{minipage}\hfill
\begin{minipage}[t]{0.48\linewidth}
	\centering
	\fbox{\parbox[c][4.0cm][c]{0.92\linewidth}{\centering
	FEA simulation image\\placeholder}}
\end{minipage}
\caption{Wheel design and FEA simulations.}
\label{fig:wheel_and_fea_placeholders}
\end{figure}

